\cleardoublepage{}
\begin{center}
    \bfseries \zihao{3} 摘~要
\end{center}

随着多模态大模型的发展，视频大模型在视频问答、时空推理与复杂场景理解等任务中展现出较强能力。然而，相较于图像场景，视频数据具有更长的时序长度与更高的视觉冗余，导致视频大模型在推理过程中面临视觉 token 数量庞大、显存占用高以及推理延迟大的问题，严重制约了其实际部署与应用。针对上述问题，本文围绕多模态视频大模型的高效推理展开研究，重点研究视觉 token 压缩技术在视频场景中的迁移与优化方法。

本文以视频大模型 VideoLLaMA2 为基础框架，引入 TokenPacker 中提出的视觉 token 压缩思想，设计了一种适用于视频场景的视觉 token 高效压缩方案。本文在不改变原有大语言模型结构与时空融合模块的前提下，将 TokenPacker 模块接入视觉编码器与时空连接模块之间，对视频帧中的空间视觉 token 进行压缩，以降低后续时空建模阶段的计算开销。针对视频特征与图像特征在时序结构上的差异，本文对原始 TokenPacker 的输入输出形式进行了适配与重构，实现了视频特征在时空维度下的稳定压缩与重组。

此外，为缓解压缩后视觉语义信息损失带来的性能下降问题，本文进一步提出基于 token 分布对齐的校准方法，通过约束压缩后 token 与原始视觉 token 在统计分布上的一致性，提高压缩特征对原模型的适配能力。同时，本文结合轻量化微调策略，对压缩模块进行进一步优化，以提升高压缩倍率场景下模型的推理性能与鲁棒性。

实验部分基于 MVBench 等视频理解数据集展开验证，并从推理准确率、视觉 token 数量、压缩倍率与推理耗时等多个维度进行分析。实验结果表明，本文方法能够在显著降低视觉 token 数量与推理开销的同时，较好地保持模型的视频理解能力；相比简单的均值池化压缩方法，本文方法在高压缩倍率场景下具有更稳定的性能表现。研究结果验证了视觉 token 压缩技术在视频大模型高效推理中的可行性，为后续视频多模态模型的轻量化部署与高效推理提供了参考。

\cleardoublepage{}
\begin{center}
    \bfseries \zihao{3} Abstract
\end{center}
